% GENERATED FILE. Edit canonical lessons, then run npm run generate:books.
% canonical-source-hash: fnv1a64:4f8154dab6123ded
% canonical-lessons: FA-C520-navazande, FA-C520-ahangsaz, FA-C520-orkestr, FA-C520-goruh-e-kor, FA-C520-albom

\chapter{Musicians}
\label{ch:fa-musicians}
\hlchaptermodality{520}

\begin{chapteropening}
\textbf{By the end of this chapter:} \emph{I can talk about musicians and their work.}

Complete the last lesson of chapter 520: I can talk about musicians and their work.
\end{chapteropening}

\section[\texorpdfstring{\fa{نوازنده} (navâzande)}{نوازنده (navâzande)}]{\fa{نوازنده} (navâzande) — a musician (who plays)}

\label{lesson:FA-C520-navazande}



\begin{quote}
Before the new one: say the Persian for rhythm, then the Persian for a musical note.
\end{quote}

\subsection*{You'll want to know: \fa{نوازنده}}
\textbf{\fa{نوازنده}} (\emph{navâzande}) — \textquotedblleft{}a musician (who plays)\textquotedblright{}.

\begin{quote}
\fa{نوازنده} \fa{جوان} \fa{در} \fa{کنسرت} \fa{ویولن} \fa{زد}.
\end{quote}

\begin{quote}
\emph{navâzande-ye javân dar konsert viyolon zad.}
\end{quote}

\begin{quote}
The young musician played the violin at the concert.
\end{quote}

\emph{navâz-} (stem of \emph{navâkhtan}) + \emph{-ande} 'one who'. A singer is \emph{khânande}.

\begin{cousinweb}
\emph{navâzande} is built on native Persian \emph{navâz-}, the present stem of \emph{navâkhtan}, 'to play music', with the ending \emph{-ande}, 'one who'. EtymDB has no entry for it.
\end{cousinweb}

\begin{grammarlens}[title={What you've built}]
One more word for telling, explaining and describing at B1.
\end{grammarlens}

\subsection*{Your turn}
\begin{itemize}
\raggedright
  \item \emph{Say it:} \emph{navâzande}
  \item \emph{Say it:} \emph{navâzande}, once more
  \item \emph{Say it:} a sentence of your own with \emph{navâzande}
  \item \emph{Recall:} say the Persian for rhythm, then the Persian for a musical note, then say \emph{navâzande} again
\end{itemize}

\begin{tcolorbox}[breakable,colback=teal!4,colframe=teal!35!black,title={Before you move on}]
What does \fa{نوازنده} mean? (\textquotedblleft{}A musician (who plays)\textquotedblright{}.) Say it once more.
\end{tcolorbox}
\section[\texorpdfstring{\fa{آهنگساز} (âhangsâz)}{آهنگساز (âhangsâz)}]{\fa{آهنگساز} (âhangsâz) — a composer}

\label{lesson:FA-C520-ahangsaz}



\begin{quote}
Before the new one: say the Persian for a musical note, then the Persian for a musician.
\end{quote}

\subsection*{You'll want to know: \fa{آهنگساز}}
\textbf{\fa{آهنگساز}} (\emph{âhangsâz}) — \textquotedblleft{}a composer\textquotedblright{}.

\begin{quote}
\fa{این} \fa{آهنگساز} \fa{موسیقی} \fa{فیلم} \fa{میسازد}.
\end{quote}

\begin{quote}
\emph{in âhangsâz musiqi-ye film misâzad.}
\end{quote}

\begin{quote}
This composer writes film music.
\end{quote}

\emph{âhang} 'tune' + \emph{sâz} 'maker' (stem of \emph{sâkhtan}).

\begin{cousinweb}
\emph{âhangsâz} is built on native Persian \emph{âhang}, which EtymDB traces to a Middle Persian word for 'pull, draught; melody' (EtymDB), and \emph{sâz}, the present stem of \emph{sâkhtan}, 'to make'.
\end{cousinweb}

\begin{grammarlens}[title={What you've built}]
One more word for telling, explaining and describing at B1.
\end{grammarlens}

\subsection*{Your turn}
\begin{itemize}
\raggedright
  \item \emph{Say it:} \emph{âhangsâz}
  \item \emph{Say it:} \emph{âhangsâz}, once more
  \item \emph{Say it:} a sentence of your own with \emph{âhangsâz}
  \item \emph{Recall:} say the Persian for a musical note, then the Persian for a musician, then say \emph{âhangsâz} again
\end{itemize}

\begin{tcolorbox}[breakable,colback=teal!4,colframe=teal!35!black,title={Before you move on}]
What does \fa{آهنگساز} mean? (\textquotedblleft{}A composer\textquotedblright{}.) Say it once more.
\end{tcolorbox}
\section[\texorpdfstring{\fa{ارکستر} (orkestr)}{ارکستر (orkestr)}]{\fa{ارکستر} (orkestr) — an orchestra}

\label{lesson:FA-C520-orkestr}



\begin{quote}
Before the new one: say the Persian for a musician, then the Persian for a composer.
\end{quote}

\subsection*{You'll want to know: \fa{ارکستر}}
\textbf{\fa{ارکستر}} (\emph{orkestr}) — \textquotedblleft{}an orchestra\textquotedblright{}.

\begin{quote}
\fa{ارکستر} \fa{شهر} \fa{هر} \fa{ماه} \fa{یک} \fa{کنسرت} \fa{دارد}.
\end{quote}

\begin{quote}
\emph{orkestr-e shahr har mâh yek konsert dârad.}
\end{quote}

\begin{quote}
The city orchestra gives a concert every month.
\end{quote}

\emph{rahbar-e orkestr} — a conductor.

\begin{cousinweb}
\emph{orkestr} is a modern loan, usually traced to French \emph{orchestre}. EtymDB lists it without a source.
\end{cousinweb}

\begin{grammarlens}[title={What you've built}]
One more word for telling, explaining and describing at B1.
\end{grammarlens}

\subsection*{Your turn}
\begin{itemize}
\raggedright
  \item \emph{Say it:} \emph{orkestr}
  \item \emph{Say it:} \emph{orkestr}, once more
  \item \emph{Say it:} a sentence of your own with \emph{orkestr}
  \item \emph{Recall:} say the Persian for a musician, then the Persian for a composer, then say \emph{orkestr} again
\end{itemize}

\begin{tcolorbox}[breakable,colback=teal!4,colframe=teal!35!black,title={Before you move on}]
What does \fa{ارکستر} mean? (\textquotedblleft{}An orchestra\textquotedblright{}.) Say it once more.
\end{tcolorbox}
\section[\texorpdfstring{\fa{گروه} \fa{کر} (goruh-e kor)}{گروه کر (goruh-e kor)}]{\fa{گروه} \fa{کر} (goruh-e kor) — a choir}

\label{lesson:FA-C520-goruh-e-kor}



\begin{quote}
Before the new one: say the Persian for a composer, then the Persian for an orchestra.
\end{quote}

\subsection*{You'll want to know: \fa{گروه} \fa{کر}}
\textbf{\fa{گروه} \fa{کر}} (\emph{goruh-e kor}) — \textquotedblleft{}a choir\textquotedblright{}.

\begin{quote}
\fa{دخترم} \fa{در} \fa{گروه} \fa{کر} \fa{مدرسه} \fa{آواز} \fa{میخواند}.
\end{quote}

\begin{quote}
\emph{dokhtaram dar goruh-e kor-e madrese âvâz mikhânad.}
\end{quote}

\begin{quote}
My daughter sings in the school choir.
\end{quote}

\emph{goruh} 'group' + \emph{kor} 'choir'. Not \emph{kar} 'deaf', spelled the same.

\begin{cousinweb}
\emph{goruh-e kor} joins inherited \emph{goruh}, 'group', which EtymDB traces to a Middle Persian word for 'group, crowd' (EtymDB), and \emph{kor}, a modern loan usually traced to French \emph{chœur}, 'choir'. EtymDB's entry for this spelling is the other word, \emph{kar}, 'deaf'.
\end{cousinweb}

\begin{grammarlens}[title={What you've built}]
One more word for telling, explaining and describing at B1.
\end{grammarlens}

\subsection*{Your turn}
\begin{itemize}
\raggedright
  \item \emph{Say it:} \emph{goruh-e kor}
  \item \emph{Say it:} \emph{goruh-e kor}, once more
  \item \emph{Say it:} a sentence of your own with \emph{goruh-e kor}
  \item \emph{Recall:} say the Persian for a composer, then the Persian for an orchestra, then say \emph{goruh-e kor} again
\end{itemize}

\begin{tcolorbox}[breakable,colback=teal!4,colframe=teal!35!black,title={Before you move on}]
What does \fa{گروه} \fa{کر} mean? (\textquotedblleft{}A choir\textquotedblright{}.) Say it once more.
\end{tcolorbox}
\section[\texorpdfstring{\fa{آلبوم} (âlbom)}{آلبوم (âlbom)}]{\fa{آلبوم} (âlbom) — an album}

\label{lesson:FA-C520-albom}



\begin{quote}
Before the new one: say the Persian for an orchestra, then the Persian for a choir.
\end{quote}

\subsection*{You'll want to know: \fa{آلبوم}}
\textbf{\fa{آلبوم}} (\emph{âlbom}) — \textquotedblleft{}an album\textquotedblright{}.

\begin{quote}
\fa{آلبوم} \fa{تازه} \fa{این} \fa{خواننده} \fa{را} \fa{خریدم}.
\end{quote}

\begin{quote}
\emph{âlbom-e tâze-ye in khânande râ kharidam.}
\end{quote}

\begin{quote}
I bought this singer's new album.
\end{quote}

Also a photo album: \emph{âlbom-e 'aks}.

\begin{cousinweb}
\emph{âlbom} is a modern loan, usually traced to French \emph{album}. EtymDB has no entry for the Persian word.
\end{cousinweb}

\begin{grammarlens}[title={What you've built}]
One more word for telling, explaining and describing at B1.
\end{grammarlens}

\subsection*{Your turn}
\begin{itemize}
\raggedright
  \item \emph{Say it:} \emph{âlbom}
  \item \emph{Say it:} \emph{âlbom}, once more
  \item \emph{Say it:} a sentence of your own with \emph{âlbom}
  \item \emph{Recall:} say the Persian for an orchestra, then the Persian for a choir, then say \emph{âlbom} again
\end{itemize}

\begin{tcolorbox}[breakable,colback=teal!4,colframe=teal!35!black,title={Before you move on}]
What does \fa{آلبوم} mean? (\textquotedblleft{}An album\textquotedblright{}.) Say it once more.
\end{tcolorbox}
